\inicializarSeccion{Metodología}
\tab La base de datos CalCOFI contiene datos oceanográficos extraídos de muestras de agua de mar, que son recolectados en CalCOFI. Estos datos incluyen parámetros como temperatura, salinidad, oxígeno disuelto, clorofila-a, nutrientes, y muchos más. 

Las siguientes variables fueron las seleccionadas, y el porqué:
\begin{itemize}
    \item \texttt{mean\_O2ml\_L} (cuantitativa): nos sirve para saber que el promedio de cuánto oxígeno contiene una zona ambiental marina contiene, refleja directamente el tipo de en el que dichas especies viven.
    \item \texttt{mean\_O2Sat} (cuantitativa): nos sirve para poder evaluar la cantidad de oxígeno que la fauna tiene a disponibilidad, con la saturación de este mismo.
    \item \texttt{mean\_T\_degC} (cuantitativa): nos sirve para saber si la temperatura cambia el comportamiento de la fauna; alterando su metabolismo, horarios cuando cazan, si van a procrear, etc.
    \item \texttt{mean\_Salnty} (cuantitativa): nos sirve para saber si checar el cambio de salinidad en las aguas permite saber sobre el desplazamiento de ciertas especies; apoyando al cambio parcial o total de un ecosistema.
    \item \texttt{mean\_ChlorA} (cuantitativa): nos sirve para conocer la cantidad de clorofila en el ecosistema, y así saber que tipo de ambientes son más probables de prosperar; estos ambientes son los que más materia orgánica poseen.
    \item \texttt{mean\_PO4uM} (cuantitativa): nos sirve para medir los fosfatos promedio en una roseta, ya que provoca un crecimiento masivo de algas, bloquea luz solar, y otros; resultando en una pérdida de biodiversidad, y graves daños ecológicos.
    \item \texttt{Depth\_zone} (categórica): Nos sirve para caracterizar a los tipos de especies que residen en ciertas zonas, conociendo que tan profunda es esta. Conocer este dato categóricamente nos permite clasificar mejor el hábitat en donde ciertas especies viven.
    \item \texttt{Water\_denisty\_class} (categórica): similar a la profundidad, conociendo la densidad del agua, podemos saltar a conclusiones acerca de cómo el ecosistema vive.
\end{itemize}

\tab Estas variables juegan un papel muy importante para el análisis oceánico. Un cambio significativo en cada una de estas variables nos da información muy valiosa sobre la salud del ecosistema. Estas nos describen más a fondo los cambios en el océano, a diferencia de las omitidas que categorizan cada uno de los ambientes, o bien terminan siendo redundantes. 
